\subsection{乘积测度作为测度积分的解释}

在本节中，T 和 \(T'\) 表示两个局部紧空间，\(\mu\) 是 T 上的一个正测度，\(\mu'\) 是 \(T'\) 上的一个正测度，而 \(\nu = \mu \otimes \mu'\) 是 \(X = T \times T'\) 上的乘积测度（Ch. III, §4, No.~1）。

对每个 \(t \in T\) ，映射 \(t' \mapsto (t, t')\) （\(T'\) 到 X 中）是连续且逆紧的。用 \(\lambda_t'\) 表示 \(\mu'\) 在此映射下的像；\(\lambda_t'\) 是 X 上的一个正测度，且若 \(f \in \mathcal{K}(X)\) ，则用 \(f_t\) 表示偏映射 \(t' \mapsto f(t, t')\) ，有

\[
\int f d \lambda_ {t} ^ {\prime} = \int f _ {t} d \mu^ {\prime},\tag{1}
\]

这也可用关系式 \(\lambda_t' = \varepsilon_t \otimes \mu'\) 来表达。

此外，映射 \(t \mapsto \lambda_t'(f)\) 是连续的且具有紧支集（Ch. III, §4, No.~1, 引理 2），因此映射 \(t \mapsto \lambda_t'\) （T 到 \(\mathcal{M}(\mathbf{X})\) 中的映射）是淡连续的（从而更是淡 \(\mu\) -可测的）；于是测度族 \(t \mapsto \lambda_t'\) 是 \(\mu\) -适宜的（§3, No.~1, 命题 2a)）。\(f\) 关于测度 \(\int \lambda_t' d\mu(t)\) 的积分按定义为

\[
\int \langle f, \lambda_ {t} ^ {\prime} \rangle d \mu (t) = \int d \mu (t) \int f _ {t} (t ^ {\prime}) d \mu^ {\prime} (t ^ {\prime}) = \int f (t, t ^ {\prime}) d \nu (t, t ^ {\prime})
\]

（Ch. III, §4, No.~1, 定理 2）；于是 \(\nu = \int \lambda_t' d\mu(t)\) 。

类似地，对每个元素 \(t^{\prime} \in T^{\prime}\) ，用 \(\lambda_{t^{\prime}}\) 表示 \(\mu\) 在 T 到 X 中的映射 \(t \mapsto (t, t^{\prime})\) 下的像。映射 \(t^{\prime} \mapsto \lambda_{t^{\prime}}\) 是 \(\mu^{\prime}\) -适宜的且淡连续的，且 \(\nu = \int \lambda_{t^{\prime}} d\mu^{\prime}(t^{\prime})\) 。我们将需要下列引理：

引理 1. --- 对每个数值函数 \(f \geqslant 0\) （定义在 \(X\) 上），

\[
\int^ {*} f _ {t} d \mu^ {\prime} = \int^ {*} f d \lambda_ {t} ^ {\prime}.\tag{2}
\]

由于 \(t' \mapsto (t, t')\) 是连续且逆紧的映射，这由 §4, No.~2 的命题 2 得出。

引理 2. --- 设 \(f\) 是 \(X\) 到一个拓扑空间中的映射。为使 \(f\) 是 \(\lambda_t'\) -可测的，必须且只须 \(f_t\) 是 \(\mu'\) -可测的。这是 §6, No.~2 命题 3 的推论。

引理 3. --- 设 f 是定义在 X 上、取值于 \(\overline{R}\) 或一个 Banach 空间中的函数。为使 f 是 \(\lambda_{t}^{\prime}\) -可积的，必须且只须 \(f_{t}\) 是 \(\mu^{\prime}\) -可积的，此时

\[
\int \mathbf {f} _ {t} d \mu^ {\prime} = \int \mathbf {f} d \lambda_ {t} ^ {\prime}.\tag{3}
\]

这由 §4, No.~4 的定理 2 得出，其中用到 \(t' \mapsto (t, t')\) 连续且逆紧这一事实。

注记. --- 引理 1、2、3 可以不用 §§4 和 6 的结果而由一个直接论证非常简单地证明。例如，若 \(f \in \mathcal{K}(\mathrm{T} \times \mathrm{T}')\) ，则关系式 (2) 按定义是明显的。若 \(f\) 在 \(\mathrm{X} = \mathrm{T} \times \mathrm{T}'\) 上下半连续，则只须注意 \(t' \mapsto f_t(t')\) 是诸函数 \(t' \mapsto g_t(t') = g(t, t')\) 的上包络，其中 \(g\) 取遍 \(\mathcal{K}(\mathrm{X})\) 中满足 \(0 \leqslant g \leqslant f\) 的函数所成之集。最后，对任意的 \(f\) ，注意：若 \(h \geqslant f\) 在 \(\mathrm{X}\) 上下半连续，则 \(t' \mapsto h(t, t')\) 在 \(\mathrm{T}'\) 上下半连续；反之，若 \(t' \mapsto u(t')\) 在 \(\mathrm{T}'\) 上下半连续且满足 \(u(t') \geqslant f(t, t')\) （对所有 \(t' \in \mathrm{T}'\) ），则函数 \(h\) ——其中 \(h(t, t') = u(t')\) ，\(h(t_1, t') = +\infty\) （当 \(t_1 \neq t'\) ）——在 \(\mathrm{X}\) 上下半连续且满足 \(h \geqslant f\) 。一旦证明了引理 1，由此即可推出集 \((\mathrm{T} - \{t\}) \times \mathrm{T}'\) 是 \(\lambda_t'\) -可忽略的，于是证明引理 2 和 3 就很容易了。

关系式 (3) 允许把它的两端都记作 \(\int\mathbf{f}(t,t^{\prime})d\mu^{\prime}(t^{\prime})\) 而不致引起混淆。对测度 \(\lambda_{t^{\prime}}=\mu\otimes\varepsilon_{t^{\prime}}\) 显然有类似的结果。

代替记号

\[
\int^ {*} f (t, t ^ {\prime}) d \nu (t, t ^ {\prime}), \quad \int^ {\bullet} f (t, t ^ {\prime}) d \nu (t, t ^ {\prime}), \quad \int \mathbf {f} (t, t ^ {\prime}) d \nu (t, t ^ {\prime}),
\]

我们将采用记号

\[
\iint^ {*} f (t, t ^ {\prime}) d \mu (t) d \mu^ {\prime} (t ^ {\prime}), \iint^ {\bullet} f (t, t ^ {\prime}) d \mu (t) d \mu^ {\prime} (t ^ {\prime}), \iint \mathbf {f} (t, t ^ {\prime}) d \mu (t) d \mu^ {\prime} (t ^ {\prime}),
\]

这与 Ch. III, §4, No.~1 中采用的记号一致。

把测度 \(\nu\) 解释为一个积分 \(\int \lambda_t' d\mu(t)\) ，将使我们可以把 §3 的结果翻译成乘积测度的语言。另一方面，测度 \(\lambda_t'\) 由 \(\{t\} \times T' = \overline{\text{pr}}_1(t)\) 承载，故这个积分定义了 \(\nu\) 的一个切片分解，它是关于投影 \(\text{pr}_1\) （即 \(T \times T'\) 到 \(T\) 上的投影）的（§6, No.~6）。在给出由此所得结果的清单之前，先给出一个有用的性质：

命题 1. --- 设 \((\mu_{\alpha})_{\alpha \in \mathrm{A}}\) （相应地，\((\mu_{\beta}^{\prime})_{\beta \in \mathrm{B}})\) ）是 T 上（相应地，\(\mathbf{T}'\) 上）正测度组成的可和族，其和记作 \(\mu\) （相应地，\(\mu'\) ）。则族 \((\mu_{\alpha} \otimes \mu_{\beta}^{\prime})_{(\alpha, \beta) \in \mathrm{A} \times \mathrm{B}}\) 在 \(\mathbf{T} \times \mathbf{T}'\) 上是可和的，且

\[
\mu \times \mu^ {\prime} = \sum_ {(\alpha , \beta) \in \mathrm{A} \times \mathrm{B}} \mu_ {\alpha} \otimes \mu_ {\beta} ^ {\prime}.\tag{4}
\]

当 A 和 B 有限时，这些性质是明显的。由此推出，若 \(A'\) （相应地，\(B'\) ）是 A（相应地，B）的一个有限子集，则

\[
\sum_ {(\alpha , \beta) \in A ^ {\prime} \times B ^ {\prime}} \mu_ {\alpha} \otimes \mu_ {\beta} ^ {\prime} \leqslant \mu \otimes \mu^ {\prime}.
\]

于是族 \((\mu_{\alpha} \otimes \mu'_{\beta})\) 是可和的。为证明 (4) 的两端相等，只须证明第二端满足乘积测度的特征性质（Ch. III, §4, No.~1, 定理 1），这由下面的计算得出。

设 \(f\) 是 \(\mathcal{H}_{+}(\mathrm{T})\) 的一个元素，\(f'\) 是 \(\mathcal{H}_{+}(\mathrm{T}')\) 的一个元素；回忆 \(f \otimes f'\) 表示函数 \((t, t') \mapsto f(t)f'(t')\) （在 \(\mathrm{T} \times \mathrm{T}'\) 上），它属于 \(\mathcal{H}_{+}(\mathrm{T} \times \mathrm{T}')\) （A, II, §7, No.~7）。于是，由乘积测度的定义，

\[
\begin{array}{r l} \sum_ {(\alpha , \beta) \in \mathrm{A} \times \mathrm{B}} \langle \mu_ {\alpha} \otimes \mu_ {\beta} ^ {\prime}, f \otimes f ^ {\prime} \rangle & = \sum_ {(\alpha , \beta) \in \mathrm{A} \times \mathrm{B}} \big (\langle \mu_ {\alpha}, f \rangle \langle \mu_ {\beta} ^ {\prime}, f ^ {\prime} \rangle \big) \\ & = \Big (\sum_ {\alpha \in \mathrm{A}} \langle \mu_ {\alpha}, f \rangle \Big) \Big (\sum_ {\beta \in \mathrm{B}} \langle \mu_ {\beta} ^ {\prime}, f ^ {\prime} \rangle \Big) \\ & = \langle \mu , f \rangle \langle \mu^ {\prime}, f ^ {\prime} \rangle \\ & = \langle \mu \otimes \mu^ {\prime}, f \otimes f ^ {\prime} \rangle . \end{array}
\]

